\newpage

\section*{Выводы и практические рекомендации}
\addcontentsline{toc}{section}{Выводы и практические рекомендации}

\subsection*{Описание данных и их типов}
 
В работе был исследован датасет \texttt{Wine Quality}, загруженный из файла \texttt{wine\_quality\_merged.csv}. Датасет содержит 6497 наблюдений и 13 признаков. Одиннадцать признаков являются количественными непрерывными признаками, измеренными по шкале отношений. Признак \texttt{quality} является категориальным порядковым признаком. Признак \texttt{type} является категориальным номинальным признаком.
 
\subsection*{Результаты проверки базовых предположений}
 
Проверка с помощью \texttt{4-Plot} показала, что исходный порядок наблюдений не является случайным. Это связано с тем, что в датасете сначала расположены объекты одного типа вина, а затем другого. Гистограммы и \texttt{Normal Probability Plot} показали, что признаки имеют многомодальные распределения в общей выборке. Для большинства признаков соответствие нормальному распределению является только приблизительным, также имеются отклонения в хвостах.
 
\subsection*{Описательные статистики и их интерпретация}
 
Для количественных признаков были рассчитаны меры центральной тенденции и меры разброса. Среднее и медиана большинства признаков имеют различные значения. Это указывает на асимметричность их распределений. Наибольший коэффициент вариации наблюдается у признаков \texttt{residual sugar}, \texttt{chlorides} и \texttt{total sulfur dioxide}. Для признака \texttt{density} коэффициент вариации наименьший.

\subsection*{Визуализации и основные результаты}
 
В работе были построены следующие визуализации: гистограммы, \texttt{Box Plot}, эмпирические функции распределения, графики \texttt{KDE}, диаграммы рассеяния, тепловая карта корреляций, матрица диаграмм рассеяния, \texttt{Box Plot} и \texttt{CDF} по типам вина. Гистограммы и \texttt{KDE} подтвердили наличие нескольких областей концентрации значений для большинства признаков. Матрица диаграмм рассеяния показала, что признаки \texttt{alcohol} и \texttt{volatile acidity} лучше всего разделяют типы вина.
 
\subsection*{Аномалии, выбросы и пропуски}
 
В датасете отсутствуют пропущенные значения. Также было обнаружено 1169 полных дублирующихся наблюдений.
 
По правилу $1{,}5 \cdot IQR$ потенциальные выбросы были обнаружены у большинства признаков. Данные наблюдения не удалялись автоматически, поскольку статистическое правило выделяет только потенциальные выбросы. Они могут быть корректными измерениями для отдельных типов вина. После исключения потенциальных выбросов значения корреляций изменились незначительно.
 
\subsection*{Гипотезы для дальнейшего анализа}
 
По результатам исследования были сформулированы следующие гипотезы:

\begin{enumerate}
    \item Между \texttt{free sulfur dioxide} и \texttt{total sulfur dioxide} существует сильная положительная связь.
 
    \item Признак \texttt{alcohol} позволяет хорошо различать типы вина.
 
    \item Признаки \texttt{volatile acidity} и \texttt{citric acid} также могут быть информативны для различения типов вина.
 
    \item Потенциальные выбросы в признаках не оказывают существенного влияния на корреляции между признаками.
 
    \item Признаки \texttt{alcohol} и \texttt{volatile acidity} являются наиболее информативными для классификации типов вина.
 
\end{enumerate}
 
\subsection*{Общий вывод о применимости моделей}
 
Исследуемый датасет пригоден для дальнейшего применения методов статистического анализа и классификации. В нём отсутствуют пропуски, группы имеют различный размер, а признаки \texttt{alcohol} и \texttt{volatile acidity} позволяют хорошо различать типы вина.
 
При построении моделей необходимо учитывать несколько особенностей данных. Во-первых, наблюдения в исходном датасете упорядочены по типу вина, поэтому для обучения моделей выборку следует перемешивать и разделять на обучающую и тестовую части с сохранением пропорций групп. Во-вторых, общие распределения признаков являются многомодальными. Поэтому их желательно анализировать отдельно по типам или использовать методы, не требующие строгого предположения о нормальности распределения. В-третьих, для методов, основанных на расстояниях или коэффициентах регрессии, может потребоваться стандартизация количественных признаков. Потенциальные выбросы следует проверять с учётом предметного смысла, а не удалять автоматически.